\documentclass[11pt]{article}

\usepackage[a4paper,margin=1in]{geometry}
\usepackage{amsmath,amssymb,amsthm,mathtools}
\usepackage{booktabs}
\usepackage{microtype}
\usepackage{hyperref}
\IfFileExists{cleveref.sty}{\usepackage[nameinlink,noabbrev]{cleveref}}{}
\providecommand{\cref}[1]{\ref{#1}}
\providecommand{\Cref}[1]{\ref{#1}}

\title{Fast Single-Point Evaluation of Univariate Polynomials\\
in the Boolean-Kernel Basis}

\author{Ali Mkhida\\
Algorizk Labs}

\date{\today}

\newtheorem{theorem}{Theorem}
\newtheorem{lemma}[theorem]{Lemma}
\newtheorem{proposition}[theorem]{Proposition}
\newtheorem{definition}[theorem]{Definition}
\newtheorem{remark}[theorem]{Remark}

\begin{document}

\maketitle

\begin{abstract}
We study single-point evaluation of large univariate polynomials represented
in a Boolean-kernel basis and compare it with conventional univariate
representations relevant to zero-knowledge proof systems. For
\(N=2^m\), the basis is indexed by \(y\in\{0,1\}^m\) and defined by
\[
K_y(X)=\prod_{i=0}^{m-1}\left(X^{2^i}+y_i\right).
\]
Its structure gives a binary evaluation tree whose elementary fold is
\[
(a,b)\longmapsto t(a+b)+b,
\qquad t=x^{2^i}.
\]
We implement and optimize this evaluator over the Goldilocks field.
For \(N=2^{20}\) on an Apple M1 processor, the final same-run
cross-basis benchmark gives a central estimate of \(1.0560\,\mathrm{ms}\)
for the optimized Boolean-kernel evaluator and \(1.4116\,\mathrm{ms}\)
for the parallel blocked monomial baseline.  This corresponds to a
\(1.337\times\) evaluation-speed advantage, or approximately \(25.2\%\)
lower wall-clock time.  A separate optimization-only run reached
\(892.57\,\mu\mathrm{s}\), illustrating the observed run-to-run variation. We describe the evaluation algorithm, its exact
arithmetic cost, the implementation optimizations responsible for the
measured speedup, and the limitations of transferring this primitive-level
gain to complete proof systems.
\end{abstract}

\section{Introduction}
\label{sec:introduction}

Polynomial arithmetic is a central computational component of modern
zero-knowledge proof systems. Provers manipulate large polynomials during
arithmetization, constraint evaluation, low-degree testing, polynomial
commitment, quotient construction, and opening procedures. Transparent
proof systems based on Reed--Solomon proximity testing, including
FRI-based constructions and STARKs, make polynomial evaluation and encoding
particularly visible components of the prover computation
\cite{stark2018}.

The representation chosen for a polynomial determines the cost of many of
these operations. A polynomial may be represented by its monomial
coefficients, by its values on a structured evaluation domain, or by
coefficients in another basis adapted to a particular algebraic operation.
These representations describe the same polynomial space but expose
different arithmetic structure and different opportunities for
parallelism.

This work isolates one basic operation:
\begin{quote}
Given a degree-\(<N\) univariate polynomial \(P\) and an arbitrary point
\(x\) in the base field, compute \(P(x)\) from the chosen representation.
\end{quote}

We deliberately study this primitive independently of a complete proof
system. This allows the effect of the representation itself to be measured
under a common field implementation, machine, input size, and benchmarking
methodology.

Let
\[
N=2^m.
\]
In addition to the monomial representation and a roots-of-unity
Lagrange/evaluation representation, we study the following family of
univariate polynomials:
\[
K_y(X)
=
\prod_{i=0}^{m-1}
\left(X^{2^i}+y_i\right),
\qquad
y=(y_0,\ldots,y_{m-1})\in\{0,1\}^m.
\]
We refer to this family as the \emph{Boolean-kernel basis}. A polynomial
represented in this basis is written
\[
P(X)
=
\sum_{y\in\{0,1\}^m}
c_y K_y(X).
\]

The basis admits a simple binary reduction rule. Fix an evaluation point
\(x\), and define
\[
t_i=x^{2^i}.
\]
At level \(i\), two coefficients \(a\) and \(b\) whose indices differ only
in bit \(i\) are replaced by
\[
F_{t_i}(a,b)
=
t_i(a+b)+b.
\]
Applying this rule recursively reduces \(N\) coefficients to the single
value \(P(x)\). The full tree contains exactly \(N-1\) folds, and all
folds at the same level are independent.

This structure makes the evaluation algorithm both linear in the number of
input coefficients and naturally parallel. Since an arbitrary coefficient
vector must in general be inspected in full, an \(O(N)\) algorithm already
has optimal asymptotic input complexity for this task. The relevant
question is therefore whether the basis gives a practically favorable
constant factor and memory-access pattern.

We answer this experimentally using a Rust implementation over the
Goldilocks field
\[
\mathbb{F}_p,
\qquad
p=2^{64}-2^{32}+1.
\]
The optimized implementation combines coarse-grained subtree parallelism,
a specialized Goldilocks reduction, a fused kernel-fold operation,
bounds-check-free inner loops, direct first-level reduction, and
worker-local reusable scratch storage.

For
\[
N=2^{20}=1{,}048{,}576,
\]
the final same-run cross-basis benchmark on the tested Apple M1 machine
gives a central estimate of
\[
\boxed{1.0560\,\mathrm{ms}}
\]
for the optimized Boolean-kernel evaluator and
\[
1.4116\,\mathrm{ms}
\]
for the blocked parallel monomial evaluator.  Thus the kernel evaluator is
\[
\boxed{1.337\times}
\]
faster in this controlled comparison, corresponding to approximately
\[
\boxed{25.2\%}
\]
less wall-clock evaluation time.  During kernel-only optimization, a
separate run reached \(892.57\,\mu\mathrm{s}\); we report that value as a
best observed implementation result but do not mix it with competitor
measurements from another run when computing cross-basis speedups.

This comparison is intentionally an operation-level result. It does
\emph{not} imply a \(1.337\times\) speedup for a complete STARK, FRI prover,
or polynomial commitment scheme. A complete protocol also contains
encoding, low-degree extension, hashing, commitment, transcript, query,
and other polynomial operations, and a change of representation may create
conversion costs elsewhere. The result instead establishes that
representation alone can materially change the cost of arbitrary-point
evaluation and motivates studying whether this basis can be maintained
across a larger fraction of a prover pipeline.

Related work on transparent proof systems already demonstrates the
importance of linear- and quasi-linear-time polynomial machinery
\cite{stark2018}. Other polynomial commitment constructions have also
highlighted the cost of moving between polynomial representations; for
example, Zeromorph explicitly addresses interactions between multilinear
polynomials and univariate commitment machinery \cite{zeromorph2023}.
Our scope is different: we experimentally isolate a univariate basis and
measure its effect on a concrete polynomial operation.

\paragraph{Contributions.}
The contributions of this work are:
\begin{enumerate}
    \item We give a direct binary-fold evaluation algorithm for the
    univariate Boolean-kernel basis and derive its exact arithmetic cost.

    \item We provide an optimized Rust implementation over the Goldilocks
    field with correctness tests and reproducible benchmarks.

    \item We compare single-point evaluation against conventional
    univariate representations under a common implementation methodology,
    including a parallel monomial baseline rather than an artificially
    serialized competitor.

    \item We provide an optimization ablation covering arithmetic reduction,
    memory organization, bounds checks, scratch-buffer reuse, subtree size,
    and unsuccessful tree-fusion strategies.

    \item We distinguish primitive-level speedups from end-to-end
    proof-system claims and identify the representation-conversion and
    integration questions that must be resolved before transferring the
    result to a complete prover.
\end{enumerate}

The remainder of the paper first defines the polynomial representations
formally, then derives the Boolean-kernel evaluation algorithm, describes
the implementation and benchmark methodology, reports the experimental
results, and discusses implications for proof-system design.

\section{Polynomial Representations}
\label{sec:representations}

This section fixes the notation used throughout the paper and defines the
three univariate representations compared experimentally.

\subsection{Ambient polynomial space}

Let \(\mathbb{F}\) be a field and let
\[
N=2^m
\]
for an integer \(m\ge 1\).  We work in the \(N\)-dimensional
\(\mathbb{F}\)-vector space
\[
\mathbb{F}[X]_{<N}
=
\{P\in\mathbb{F}[X] : \deg P < N\}.
\]
Every representation considered below encodes exactly one polynomial in
this space using \(N\) field elements.

For an integer
\[
e=\sum_{i=0}^{m-1} e_i 2^i,
\qquad e_i\in\{0,1\},
\]
we write
\[
\operatorname{bits}(e)=(e_0,\ldots,e_{m-1}).
\]
For Boolean vectors \(u,v\in\{0,1\}^m\), we write
\[
u\subseteq v
\]
when \(u_i\le v_i\) for every \(i\).  Equivalently, the support of \(u\)
is a subset of the support of \(v\).

\subsection{Monomial representation}

The standard monomial basis of \(\mathbb{F}[X]_{<N}\) is
\[
1,X,X^2,\ldots,X^{N-1}.
\]
A polynomial is represented by coefficients
\[
a_0,\ldots,a_{N-1}\in\mathbb{F}
\]
such that
\[
P(X)=\sum_{e=0}^{N-1}a_eX^e.
\]

For a point \(x\in\mathbb{F}\), Horner evaluation computes
\[
P(x)
=
a_0+x\bigl(a_1+x(a_2+\cdots+x a_{N-1})\bigr),
\]
or equivalently by processing the coefficients in reverse order with the
recurrence
\[
r_{j+1}=xr_j+a_{N-1-j},
\qquad r_0=0.
\]
The serial dependency chain makes ordinary Horner evaluation inherently
sequential, although a long coefficient vector can be partitioned into
blocks and the blocks evaluated in parallel.  Our experimental comparison
therefore includes both serial Horner evaluation and a blocked parallel
monomial baseline.

\subsection{Roots-of-unity Lagrange representation}

Assume that \(\mathbb{F}\) contains a primitive \(N\)-th root of unity
\(\omega\).  Define the multiplicative evaluation domain
\[
H
=
\{1,\omega,\omega^2,\ldots,\omega^{N-1}\}.
\]
A polynomial \(P\in\mathbb{F}[X]_{<N}\) can be represented by its values
\[
v_j=P(\omega^j),
\qquad 0\le j<N.
\]

The associated Lagrange basis consists of polynomials
\[
L_j(X)
=
\prod_{\substack{0\le k<N\\k\ne j}}
\frac{X-\omega^k}{\omega^j-\omega^k},
\]
which satisfy
\[
L_j(\omega^k)=
\begin{cases}
1,&j=k,\\
0,&j\ne k.
\end{cases}
\]
Hence
\[
P(X)=\sum_{j=0}^{N-1}v_jL_j(X).
\]

For a roots-of-unity domain one may also use
\[
L_j(X)
=
\frac{X^N-1}{N(X-\omega^j)}\,\omega^j,
\]
provided \(X\ne\omega^j\).  This identity is useful for arbitrary-point
evaluation because all denominators \(x-\omega^j\) can be inverted with a
single batch inversion.

This representation is especially natural when a proof system already
stores polynomial values on an FFT/NTT domain.  Our benchmark measures
arbitrary-point evaluation from this representation rather than the
separate cost of constructing the representation.

\subsection{Boolean-kernel representation}

For each Boolean vector
\[
y=(y_0,\ldots,y_{m-1})\in\{0,1\}^m,
\]
define
\[
K_y(X)
=
\prod_{i=0}^{m-1}
\left(X^{2^i}+y_i\right).
\]
There are exactly \(2^m=N\) such polynomials.

A Boolean-kernel coefficient vector is a family
\[
(c_y)_{y\in\{0,1\}^m}
\]
representing the polynomial
\[
P(X)
=
\sum_{y\in\{0,1\}^m}
c_yK_y(X).
\]

Although every \(K_y\) has degree at most \(N-1\), and in fact every
\(K_y\) contains the top monomial \(X^{N-1}\), the family is nevertheless
a basis.  Its coefficient matrix is a Boolean zeta matrix after a reversal
of the monomial index.

\begin{lemma}[Monomial coefficients of \(K_y\)]
\label{lem:kernel-monomial-coeff}
Let
\[
e=\sum_{i=0}^{m-1}e_i2^i
\]
with \(e_i\in\{0,1\}\), and define the Boolean complement vector
\[
\bar e=(1-e_0,\ldots,1-e_{m-1}).
\]
Then
\[
[X^e]K_y(X)
=
\begin{cases}
1,&\bar e\subseteq y,\\
0,&\text{otherwise}.
\end{cases}
\]
\end{lemma}

\begin{proof}
Expanding
\[
K_y(X)=\prod_{i=0}^{m-1}(X^{2^i}+y_i),
\]
the monomial \(X^e\) can arise in exactly one way because the powers
\(2^0,\ldots,2^{m-1}\) give the unique binary expansion of \(e\).

For each index \(i\) with \(e_i=1\), the term \(X^{2^i}\) must be selected.
For each index \(i\) with \(e_i=0\), the constant term \(y_i\) must be
selected.  Therefore
\[
[X^e]K_y(X)
=
\prod_{i:e_i=0}y_i.
\]
Since every \(y_i\) is either \(0\) or \(1\), this product equals \(1\)
exactly when \(y_i=1\) for every \(i\) such that \(e_i=0\).  This condition
is precisely
\[
\bar e\subseteq y.
\]
\end{proof}

\begin{theorem}[Boolean-kernel basis]
\label{thm:kernel-basis}
The family
\[
\mathcal{K}_m
=
\{K_y(X):y\in\{0,1\}^m\}
\]
is a basis of \(\mathbb{F}[X]_{<N}\).
\end{theorem}

\begin{proof}
There are exactly \(N\) polynomials in \(\mathcal{K}_m\), so it is enough
to prove linear independence.

Index the monomial rows not directly by \(e\), but by the Boolean vector
\[
z=\bar e.
\]
By \cref{lem:kernel-monomial-coeff}, the coefficient matrix \(A\) of the
family \(\mathcal{K}_m\) in the monomial basis is
\[
A_{z,y}
=
\mathbf{1}[z\subseteq y].
\]
This is the zeta matrix of the Boolean subset lattice.

Order the Boolean vectors by nondecreasing Hamming weight, breaking ties
arbitrarily.  Whenever \(z\subseteq y\), either \(z=y\) or
\[
|z|<|y|.
\]
Consequently, under this ordering, \(A\) is upper triangular with diagonal
entries
\[
A_{y,y}=1.
\]
Hence
\[
\det A=1\ne0
\]
in every field, so \(A\) is invertible.  Therefore the polynomials
\(K_y\) are linearly independent and form a basis of
\(\mathbb{F}[X]_{<N}\).
\end{proof}

\subsection{Relation to the Boolean zeta transform}

The proof above also gives an explicit relationship between monomial and
Boolean-kernel coefficients.

Let
\[
P(X)=\sum_{e=0}^{N-1}a_eX^e
=
\sum_{y\in\{0,1\}^m}c_yK_y(X).
\]
For the exponent whose Boolean complement is \(z\), that is
\[
e=(N-1)\mathbin{\mathrm{xor}} z,
\]
\cref{lem:kernel-monomial-coeff} gives
\[
a_e
=
\sum_{y:\,z\subseteq y}c_y.
\]
Thus, after complementing the monomial index, the map from kernel
coefficients to monomial coefficients is the Boolean superset-zeta
transform.

The inverse map is Boolean Möbius inversion:
\[
c_y
=
\sum_{z:\,y\subseteq z}
(-1)^{|z|-|y|}
a_{(N-1)\mathbin{\mathrm{xor}} z}.
\]
Algorithmically, both transforms can be implemented in
\[
O(N\log N)
\]
field operations by the usual butterfly-style zeta/Möbius transform.

This conversion cost is deliberately separated from the single-point
evaluation benchmarks.  In each benchmark the polynomial is assumed to
already be available in the representation being measured.  Conversion
costs become important when considering an end-to-end proof-system
pipeline and are discussed separately later in the paper.

\subsection{A small example}

Take \(m=2\), so \(N=4\).  The four kernel polynomials are
\[
\begin{aligned}
K_{00}(X)&=X(X^2)=X^3,\\
K_{01}(X)&=(X+1)X^2=X^3+X^2,\\
K_{10}(X)&=X(X^2+1)=X^3+X,\\
K_{11}(X)&=(X+1)(X^2+1)=X^3+X^2+X+1.
\end{aligned}
\]
Here the first Boolean coordinate corresponds to the factor
\(X^{2^0}+y_0\), and the second to \(X^{2^1}+y_1\).

The coefficient pattern is triangular after complementing the monomial
index.  In particular, the constant monomial occurs only in \(K_{11}\),
the monomial \(X\) occurs in \(K_{01}\) and \(K_{11}\), the monomial
\(X^2\) occurs in \(K_{10}\) and \(K_{11}\), and \(X^3\) occurs in all
four kernel polynomials.

This subset structure is the algebraic reason that the representation is
invertible.  In the next section we show that the same product structure
also produces a particularly simple binary evaluation algorithm.

\section{Boolean-Kernel Evaluation}
\label{sec:evaluation}

We now derive the single-point evaluation algorithm for the
Boolean-kernel representation introduced in \cref{sec:representations}.

Let
\[
P(X)
=
\sum_{y\in\{0,1\}^m}
c_y
\prod_{i=0}^{m-1}
\left(X^{2^i}+y_i\right),
\]
and fix an evaluation point
\[
x\in\mathbb{F}.
\]
Define
\[
t_i=x^{2^i},
\qquad
0\le i<m.
\]

The key observation is that two basis polynomials whose indices differ
only in one Boolean coordinate can be combined locally.

\subsection{The elementary fold identity}

Fix a level \(i\).  Let
\[
u=(u_0,\ldots,u_{i-1},u_{i+1},\ldots,u_{m-1})
\in\{0,1\}^{m-1}
\]
denote the common values of all Boolean coordinates other than \(i\).
Let
\[
K_{u,0}(X)
\]
and
\[
K_{u,1}(X)
\]
denote the two kernel polynomials whose \(i\)-th coordinate is respectively
\(0\) and \(1\).

By definition,
\[
K_{u,0}(X)
=
X^{2^i}
\prod_{j\ne i}
\left(X^{2^j}+u_j\right),
\]
while
\[
K_{u,1}(X)
=
\left(X^{2^i}+1\right)
\prod_{j\ne i}
\left(X^{2^j}+u_j\right).
\]

Suppose the corresponding coefficients are
\[
a,b\in\mathbb{F}.
\]
Their contribution to \(P(x)\) is
\[
aK_{u,0}(x)+bK_{u,1}(x).
\]
Factoring the common part gives
\[
\begin{aligned}
aK_{u,0}(x)+bK_{u,1}(x)
&=
\left(
a\,x^{2^i}
+
b(x^{2^i}+1)
\right)
\prod_{j\ne i}
\left(x^{2^j}+u_j\right)
\\
&=
\left(
x^{2^i}(a+b)+b
\right)
\prod_{j\ne i}
\left(x^{2^j}+u_j\right).
\end{aligned}
\]

Therefore, after eliminating Boolean coordinate \(i\), the two
coefficients \(a\) and \(b\) are replaced by the single coefficient
\[
\boxed{
F_{t_i}(a,b)
=
t_i(a+b)+b
}
\]
with
\[
t_i=x^{2^i}.
\]

\begin{lemma}[Pair-fold identity]
\label{lem:pair-fold}
Let \(a,b\in\mathbb{F}\), let \(i\in\{0,\ldots,m-1\}\), and let \(u\)
fix all Boolean coordinates other than \(i\).  Then
\[
aK_{u,0}(x)+bK_{u,1}(x)
=
F_{t_i}(a,b)
\prod_{j\ne i}
\left(x^{2^j}+u_j\right),
\]
where
\[
F_{t_i}(a,b)=t_i(a+b)+b.
\]
\end{lemma}

\begin{proof}
The identity follows directly from the factorization above.
\end{proof}

\subsection{Recursive reduction}

The pair-fold identity can be applied recursively over all Boolean
coordinates.  We use the coefficient ordering in which the integer index
\[
y=\sum_{i=0}^{m-1}y_i2^i
\]
stores the Boolean vector
\[
(y_0,\ldots,y_{m-1}).
\]
Hence the least significant bit corresponds to the first folding level.

Let
\[
v^{(0)}_j=c_j,
\qquad
0\le j<N.
\]
After level \(i\), the active vector has length
\[
N_i=\frac{N}{2^i}.
\]

For
\[
0\le j<\frac{N_i}{2},
\]
define
\[
v^{(i+1)}_j
=
F_{t_i}
\left(
v^{(i)}_{2j},
v^{(i)}_{2j+1}
\right).
\]
Equivalently,
\[
v^{(i+1)}_j
=
t_i
\left(
v^{(i)}_{2j}
+
v^{(i)}_{2j+1}
\right)
+
v^{(i)}_{2j+1}.
\]

After \(m\) levels only one field element remains.

\begin{theorem}[Correctness of Boolean-kernel evaluation]
\label{thm:kernel-eval-correct}
Let
\[
P(X)
=
\sum_{y\in\{0,1\}^m}
c_yK_y(X).
\]
Construct the vectors
\[
v^{(0)},v^{(1)},\ldots,v^{(m)}
\]
by the recurrence above using
\[
t_i=x^{2^i}.
\]
Then
\[
v^{(m)}_0=P(x).
\]
\end{theorem}

\begin{proof}
We prove a stronger invariant.

After completing levels
\[
0,1,\ldots,i-1,
\]
the vector \(v^{(i)}\) is indexed by the remaining Boolean coordinates
\[
(y_i,\ldots,y_{m-1}).
\]
For every remaining index
\[
z=(z_i,\ldots,z_{m-1})\in\{0,1\}^{m-i},
\]
the value \(v^{(i)}_z\) is the coefficient of
\[
\prod_{j=i}^{m-1}
\left(X^{2^j}+z_j\right)
\]
after the already-eliminated coordinates
\[
0,\ldots,i-1
\]
have been evaluated at \(x\).

The invariant is true for \(i=0\) by the original representation of
\(P\).

Assume it holds for some \(i<m\).  Pair entries whose remaining indices
agree in every coordinate except \(y_i\).  By
\cref{lem:pair-fold}, replacing such a pair
\[
(a,b)
\]
by
\[
F_{t_i}(a,b)
\]
exactly evaluates the factor associated with coordinate \(i\), while
leaving the factors for coordinates
\[
i+1,\ldots,m-1
\]
unchanged.  Therefore the invariant holds for \(i+1\).

After \(m\) levels no Boolean coordinates remain, so the unique surviving
value is exactly
\[
P(x).
\]
\end{proof}

\subsection{Exact arithmetic complexity}

At level \(i\), the active vector has length
\[
\frac{N}{2^i},
\]
so the number of folds is
\[
\frac{N}{2^{i+1}}.
\]
Summing over all \(m\) levels gives
\[
\sum_{i=0}^{m-1}
\frac{N}{2^{i+1}}
=
N
\sum_{k=1}^{m}
2^{-k}
=
N\left(1-\frac{1}{2^m}\right)
=
N-1.
\]

Therefore the evaluation tree contains exactly
\[
\boxed{N-1}
\]
folds.

Each fold
\[
F_t(a,b)=t(a+b)+b
\]
contains, at the field-operation level,

\[
1
\]
multiplication and
\[
2
\]
additions.

Hence the tree itself requires exactly
\[
\boxed{N-1}
\]
field multiplications and
\[
\boxed{2(N-1)}
\]
field additions.

The level parameters are generated by
\[
t_0=x,
\qquad
t_{i+1}=t_i^2.
\]
Thus producing all
\[
t_0,\ldots,t_{m-1}
\]
requires
\[
\boxed{m-1}
\]
field squarings.

We summarize the operation count in the following proposition.

\begin{proposition}[Exact operation count]
\label{prop:operation-count}
For \(N=2^m\), Boolean-kernel single-point evaluation requires
\[
N-1
\]
kernel folds.  Using the direct fold formula, this corresponds to
\[
N-1
\]
field multiplications,
\[
2(N-1)
\]
field additions, and
\[
m-1
\]
field squarings for generation of the level parameters \(t_i\).
\end{proposition}

\subsection{Parallel structure}

The dependency graph is a complete binary reduction tree.  At each fixed
level \(i\), all folds operate on disjoint coefficient pairs and are
therefore independent.

The number of available independent folds at level \(i\) is
\[
\frac{N}{2^{i+1}}.
\]
Thus the lower levels expose abundant data parallelism, while the final
levels contain progressively less work.

A naive parallel implementation may synchronize globally after every tree
level.  For large \(N\), however, such synchronization can dominate the
shrinking amount of useful work near the top of the tree.  The optimized
implementation studied later therefore groups multiple consecutive levels
inside independent power-of-two subtrees and parallelizes only across
those larger subtrees.

This changes the scheduling strategy but not the arithmetic graph or the
number of folds.

\subsection{Input complexity}

The evaluation algorithm uses
\[
O(N)
\]
field operations and reads an arbitrary vector of \(N\) coefficients.

For the unrestricted problem considered here, this linear dependence on
the input length is unavoidable in the worst case.  To see this, suppose
an algorithm does not inspect some coefficient \(c_y\).  For any evaluation
point \(x\) such that
\[
K_y(x)\ne0,
\]
changing only \(c_y\) changes
\[
P(x)
\]
while leaving all inspected coefficients unchanged.  Hence an algorithm
that is correct on arbitrary coefficient vectors must in general obtain
information about every coefficient.

Accordingly, the principal optimization target is not the asymptotic
complexity but the constant factor: field reduction, memory traffic,
allocation, bounds checking, scheduling, and the organization of the fold
tree.

\subsection{Example for \(N=8\)}

Let
\[
m=3,
\qquad
N=8.
\]
The input coefficient vector is
\[
(c_0,c_1,c_2,c_3,c_4,c_5,c_6,c_7).
\]
Set
\[
t_0=x,
\qquad
t_1=x^2,
\qquad
t_2=x^4.
\]

The first level computes
\[
\begin{aligned}
u_0&=F_{t_0}(c_0,c_1),\\
u_1&=F_{t_0}(c_2,c_3),\\
u_2&=F_{t_0}(c_4,c_5),\\
u_3&=F_{t_0}(c_6,c_7).
\end{aligned}
\]

The second level computes
\[
\begin{aligned}
w_0&=F_{t_1}(u_0,u_1),\\
w_1&=F_{t_1}(u_2,u_3).
\end{aligned}
\]

Finally,
\[
P(x)=F_{t_2}(w_0,w_1).
\]

The tree uses
\[
4+2+1=7=N-1
\]
folds, illustrating the general count of
\cref{prop:operation-count}.

\section{Implementation and Optimization}
\label{sec:implementation}

This section describes the Rust implementation used in the experiments and
the sequence of optimizations that led to the final evaluator.  The goal is
not merely to present the fastest implementation, but to separate
algorithmic structure from low-level engineering effects and to document
which optimization ideas did and did not improve performance on the tested
Apple M1 platform.

\subsection{Implementation setting}

The implementation is written in Rust and uses Rayon for CPU parallelism.
All basis implementations in a given benchmark use the same finite-field
backend.

The experiments in this paper use the Goldilocks field
\[
\mathbb{F}_p,
\qquad
p=2^{64}-2^{32}+1.
\]
Field elements are stored canonically in a single \(64\)-bit machine word.

The optimized Boolean-kernel evaluator operates on coefficient vectors of
length
\[
N=2^m.
\]
The mathematical fold is
\[
F_t(a,b)=t(a+b)+b.
\]
The implementation therefore has two distinct optimization layers:

\begin{enumerate}
    \item optimize the field-level realization of \(F_t\);
    \item organize the binary tree so that memory traffic, allocation, and
    synchronization are minimized.
\end{enumerate}

\subsection{Baseline serial implementation}

The direct serial implementation copies the coefficient vector into a work
buffer and performs the tree level by level.

If the active prefix has length \(L\), one level computes
\[
\frac{L}{2}
\]
folds
\[
\texttt{buf[j]}
\leftarrow
F_t(\texttt{buf[2j]},\texttt{buf[2j+1]}),
\]
then replaces \(L\) by \(L/2\).

This implementation closely follows the recurrence of
\cref{sec:evaluation} and is used as a correctness reference.

\subsection{Coarse-grained subtree parallelism}

A first parallel implementation may assign all independent folds at a tree
level to Rayon and synchronize before proceeding to the next level.  This
exposes maximal parallelism at the bottom of the tree, but it also creates a
global synchronization point at every level.

Instead, the optimized implementation decomposes the coefficient vector
into contiguous power-of-two blocks of size
\[
B.
\]
Each block represents a complete local subtree.  Rayon distributes those
subtrees across workers, and each worker reduces an entire subtree locally
before returning its root.

The remaining roots form a much smaller top tree, which is finished
serially.

For
\[
N=2^{20},
\qquad
B=4096=2^{12},
\]
the input is split into
\[
\frac{2^{20}}{2^{12}}=256
\]
independent local subtrees.

This organization removes the need for a global barrier after each of the
first twelve tree levels.

\subsection{Fast Goldilocks multiplication}

A generic implementation may multiply two canonical \(64\)-bit residues
into a \(128\)-bit product and reduce the result modulo
\[
p=2^{64}-2^{32}+1.
\]

The Goldilocks modulus permits a more specialized reduction.  Write
\[
x= \ell + 2^{64}h,
\]
where
\[
0\le \ell,h<2^{64}.
\]
Let
\[
\varepsilon = 2^{32}-1.
\]
Since
\[
2^{64}\equiv \varepsilon \pmod p,
\]
we have
\[
x\equiv \ell+h\varepsilon\pmod p.
\]

Split
\[
h=h_{\mathrm{lo}}+2^{32}h_{\mathrm{hi}},
\]
with
\[
0\le h_{\mathrm{lo}},h_{\mathrm{hi}}<2^{32}.
\]
Using
\[
2^{32}\varepsilon
=
2^{64}-2^{32}
\equiv -1
\pmod p,
\]
we obtain
\[
h\varepsilon
=
h_{\mathrm{lo}}\varepsilon
+
2^{32}h_{\mathrm{hi}}\varepsilon
\equiv
h_{\mathrm{lo}}\varepsilon-h_{\mathrm{hi}}
\pmod p.
\]
Therefore
\[
\boxed{
x
\equiv
\ell-h_{\mathrm{hi}}
+
h_{\mathrm{lo}}(2^{32}-1)
\pmod p.
}
\]

This form avoids the repeated high-half multiplications generated by the
more generic reduction routine.

On the tested Apple M1 system, the multiplication microbenchmark improved
from approximately
\[
2.2372\ \mathrm{ns}
\]
to
\[
2.0559\ \mathrm{ns},
\]
an improvement of roughly
\[
8\%.
\]

\subsection{Fused kernel fold}

The direct fold
\[
F_t(a,b)=t(a+b)+b
\]
could be implemented as one modular addition, one modular multiplication,
and a second modular addition.

The optimized implementation instead fuses the final \(+b\) into the
\(128\)-bit product before Goldilocks reduction.

Let
\[
s=(a+b)\bmod p.
\]
The implementation forms the full product
\[
t s
\]
as a \(128\)-bit value, adds \(b\) to that \(128\)-bit value with carry,
and only then applies the fast Goldilocks reduction.

This preserves the mathematical value
\[
t(a+b)+b
\pmod p
\]
while avoiding a second complete modular-addition sequence after the
multiplication.

The isolated fold microbenchmark improved from approximately
\[
4.0591\ \mathrm{ns}
\]
to
\[
3.3585\ \mathrm{ns},
\]
corresponding to an improvement of about
\[
17.3\%.
\]

\subsection{Bounds-check elimination}

After arithmetic optimization, the inner tree loop became short enough
that Rust bounds checks were visible in the generated assembly.

The safe indexed loop
\[
\texttt{buf[2j]},\qquad
\texttt{buf[2j+1]},\qquad
\texttt{buf[j]}
\]
generated two conditional bounds-check branches per fold iteration on the
tested compiler target.

The optimized implementation uses a small audited helper based on raw
pointers.  For an active length \(L\), the loop is executed only for
\[
0\le j<\frac{L}{2}.
\]
Therefore
\[
2j+1<L
\]
and
\[
j<\frac{L}{2}\le L.
\]
The output location of iteration \(j\) also cannot overwrite a future input
location, because for every future index \(k>j\),
\[
2k>j.
\]

Thus the in-place access pattern is valid, and the unchecked helper removes
the per-iteration bounds checks while preserving the same sequential
semantics.

For a \(4096\)-element level, the measured runtime improved from
\[
6.9671\ \mu\mathrm{s}
\]
to
\[
6.3575\ \mu\mathrm{s},
\]
an improvement of approximately
\[
8.75\%.
\]

\subsection{Direct first-level reduction}

The first subtree implementation began each local job by copying all
\(B\) input coefficients into a private mutable buffer:
\[
B
\longrightarrow
B.
\]
It then applied the first tree level:
\[
B
\longrightarrow
\frac{B}{2}.
\]

The optimized implementation eliminates this full copy.  It reads pairs
directly from the immutable input slice and writes the first-level outputs
into a newly allocated half-size buffer:
\[
\boxed{
B
\longrightarrow
\frac{B}{2}.
}
\]

For
\[
B=4096,
\]
this changes the local scratch allocation from \(4096\) field elements to
\(2048\) field elements and eliminates the initial full-block copy.

At the first-stage microbenchmark level, the measured runtime improved from
\[
6.9031\ \mu\mathrm{s}
\]
to
\[
6.1571\ \mu\mathrm{s},
\]
an improvement of approximately
\[
10.8\%.
\]

At the full-evaluator level, this optimization reduced the best measured
runtime from approximately
\[
1.0179\ \mathrm{ms}
\]
to
\[
0.93533\ \mathrm{ms}.
\]

\subsection{Worker-local scratch reuse}

Even after direct first-level reduction, creating a new half-size vector
for every subtree still incurs repeated allocation overhead.

The final implementation uses Rayon worker-local state.  Each worker owns
a reusable scratch vector and processes multiple subtrees with that same
allocation.

This differs from an earlier unsuccessful global-buffer strategy:
the source coefficient array remains immutable and all workers still read
their assigned input subtrees directly in parallel.  Only the private
scratch allocation is reused.

In the worker-scratch optimization experiment, the same-run baseline required
\[
945.77\ \mu\mathrm{s},
\]
whereas worker-local scratch reuse required
\[
892.57\ \mu\mathrm{s}.
\]
This is the best observed kernel-only optimization run; the final
cross-basis benchmark is reported separately in \cref{sec:results}.
This corresponds to an improvement of approximately
\[
5.6\%
\]
in that run.

The resulting throughput was approximately
\[
1.175\ \mathrm{Gelem/s}.
\]

\subsection{Subtree-size tuning}

The subtree size \(B\) trades scheduling overhead against local work size
and available parallelism.

After the arithmetic and memory optimizations above, we re-ran the block
size sweep because the optimum found for an earlier, slower implementation
need not remain optimal.

A broad sweep over
\[
B\in
\{1024,2048,4096,8192,16384\}
\]
identified
\[
2048,\ 4096,\ 8192
\]
as the competitive region.

A longer controlled comparison gave:

\begin{center}
\begin{tabular}{lr}
\toprule
Block size \(B\) & Median runtime \\
\midrule
2048 & \(1.0151\ \mathrm{ms}\) \\
4096 & \(920.45\ \mu\mathrm{s}\) \\
8192 & \(917.46\ \mu\mathrm{s}\) \\
\bottomrule
\end{tabular}
\end{center}

The difference between \(4096\) and \(8192\) was only about
\[
0.32\%,
\]
which is inside the observed run-to-run noise.  We therefore retain
\[
\boxed{B=4096}
\]
as the conservative default.

\subsection{Optimization progression}

The approximate progression of the best subtree-parallel implementation at
\[
N=2^{20}
\]
was
\[
1.2530\ \mathrm{ms}
\rightarrow
1.2007\ \mathrm{ms}
\rightarrow
1.1619\ \mathrm{ms}
\rightarrow
1.0179\ \mathrm{ms}
\rightarrow
0.93533\ \mathrm{ms}
\rightarrow
\boxed{0.89257\ \mathrm{ms}}.
\]

Relative to the earlier \(1.2530\)-ms subtree baseline, the final observed
runtime is lower by approximately
\[
28.8\%.
\]

The sequence is important: no single optimization accounts for the full
gain.  Arithmetic, compiler-generated checks, memory traffic, and allocation
all become visible as preceding bottlenecks are removed.

\subsection{Rejected optimization attempts}

Several plausible optimizations did not improve performance and were not
included in the final evaluator.

\paragraph{Global shared-buffer reuse.}
A version that copied the complete coefficient vector into one shared
mutable buffer and then reduced disjoint chunks avoided per-subtree
allocation, but the initial global copy serialized too much memory traffic.
It was slower than the local-subtree design.

\paragraph{Manual four-way loop unrolling.}
Explicitly computing four independent folds per source loop iteration did
not improve the generated scalar code enough to offset the additional
register pressure and code size.

\paragraph{Alternative modular addition.}
A manually derived Goldilocks addition sequence was compared with the
compiler-generated implementation.  The compiler's existing sequence was
already shorter, so the alternative was rejected before integration.

\paragraph{Three-level tree fusion.}
A local \(8\to4\to2\to1\) fused tree reduces intermediate memory traffic,
but it also increases dependency depth and register pressure.  The fused
version was slower than the ordinary level-by-level local reduction.

\paragraph{Two-level tree fusion.}
We separately tested a smaller
\[
4\to2\to1
\]
fusion.  The current two-level path required approximately
\[
9.7019\ \mu\mathrm{s},
\]
while the fused version required
\[
9.8528\ \mu\mathrm{s}.
\]
Thus the fused variant was approximately
\[
1.6\%
\]
slower in this benchmark.

These negative results suggest that on the tested Apple M1 scalar path,
reducing memory traffic is not automatically beneficial when it lengthens
the dependency chain around the already optimized field fold.

\subsection{Final implementation}

The final evaluator used for the headline benchmark combines:

\begin{enumerate}
    \item coarse-grained subtree parallelism;
    \item block size \(B=4096\);
    \item specialized Goldilocks multiplication reduction;
    \item fused kernel-fold reduction;
    \item unchecked inner-level indexing with explicit safety invariants;
    \item direct first-level reduction into a half-size scratch buffer;
    \item Rayon worker-local scratch reuse;
    \item serial completion of the small top tree.
\end{enumerate}

All of these optimizations preserve the exact fold graph derived in
\cref{sec:evaluation}.  They change only how the arithmetic is implemented
and scheduled.

\section{Benchmark Methodology}
\label{sec:methodology}

The purpose of the benchmark methodology is to isolate the effect of
polynomial representation and implementation strategy while avoiding
artificial disadvantages for competing representations. All reported
timings are wall-clock measurements of actual Rust implementations.

\subsection{Benchmark platform}

The headline experiments were run on an Apple M1 system. The repository
records the exact compiler, operating-system, processor, memory, Git
revision, and thread-environment metadata associated with the benchmark in
\texttt{paper/benchmark-environment.txt}.

This machine-generated metadata file is committed with the paper so that
the software environment associated with the measurements can be inspected
independently of the prose description.

The benchmark is compiled in Rust release mode. The release profile uses
one code-generation unit, fat link-time optimization, and abort-on-panic
semantics:
\[
\texttt{codegen-units}=1,
\qquad
\texttt{lto}=\texttt{fat},
\qquad
\texttt{panic}=\texttt{abort}.
\]

Where assembly inspection was required during optimization, compilation was
performed for the native target CPU. The performance results reported in
the benchmark tables are obtained from the release benchmark binaries
described by the repository configuration.

\subsection{Field}

All compared polynomial representations use the same Goldilocks field
implementation:
\[
\mathbb{F}_p,
\qquad
p=2^{64}-2^{32}+1.
\]

Using a common field backend is essential. Comparing one basis with a
specialized field implementation against another basis with generic modular
arithmetic would measure field engineering rather than the representation
itself.

The final Boolean-kernel implementation uses the optimized Goldilocks
arithmetic described in \cref{sec:implementation}. Final cross-basis
numbers must therefore be collected after the same arithmetic backend is
active for every representation participating in the comparison.

\subsection{Input sizes}

The principal scaling experiment uses
\[
N\in
\{
2^{10},
2^{12},
2^{14},
2^{16},
2^{18},
2^{20}
\}.
\]

The headline comparison focuses on
\[
N=2^{20}=1{,}048{,}576.
\]

Power-of-two sizes are natural for the Boolean-kernel tree and also permit
a genuine power-of-two roots-of-unity Lagrange domain in the Goldilocks
field.

\subsection{Input generation}

For reproducibility, the monomial coefficient vector used by the main
Criterion benchmark is deterministic. For
\[
0\le i<N,
\]
the implementation constructs
\[
a_i=17i+3
\]
as a canonical field element.

The Boolean-kernel coefficient vector is obtained from the same polynomial
using the monomial-to-kernel conversion implemented in the repository.
That conversion is performed outside the timed evaluation loop.

The evaluation point used by the main benchmark is
\[
x=19.
\]

These deterministic values are chosen for reproducibility rather than for
special algebraic structure. Correctness tests separately exercise the
representation identities and optimized field arithmetic.

\subsection{Measured operation}

The principal operation is:

\begin{quote}
Given a polynomial already stored in a specified representation and an
arbitrary field point \(x\), compute \(P(x)\).
\end{quote}

The timed region therefore includes the work required to evaluate from that
representation but excludes conversion from another representation.

This distinction is intentional. The benchmark asks how expensive
single-point evaluation is once a prover already owns the polynomial in the
representation being measured. End-to-end systems may need conversion, and
those costs must be accounted for separately when evaluating a complete
prover pipeline.

\subsection{Monomial baselines}

Two monomial baselines are retained.

The first is ordinary serial Horner evaluation. It is useful because it is
the standard direct method and exposes the dependency chain inherent in
that formulation.

The second is a blocked parallel evaluator. If the block size is \(B\),
write
\[
P(X)
=
\sum_j X^{jB}P_j(X),
\]
where each block polynomial \(P_j\) has degree less than \(B\).
The values \(P_j(x)\) can be evaluated independently and combined
afterward.

The paper uses this blocked implementation as the fair parallel monomial
baseline. We do not claim a kernel speedup by comparing a parallel kernel
implementation only against an artificially serialized competitor.

For the headline experiments the monomial block size is
\[
B=4096.
\]

\subsection{Boolean-kernel configuration}

The final Boolean-kernel evaluator uses the implementation described in
\cref{sec:implementation} with local subtree size
\[
B=4096.
\]

The block-size study showed no reproducible advantage large enough to
justify replacing this setting with \(8192\), so \(4096\) is retained as
the conservative default.

Rayon provides the CPU worker pool. Unless a benchmark explicitly studies
thread scaling, the default Rayon worker count for the machine is used.
Thread-scaling experiments, when reported, must state
\texttt{RAYON\_NUM\_THREADS} explicitly.

\subsection{Criterion configuration}

The principal Criterion benchmark uses:

\begin{itemize}
    \item sample size: \(50\);
    \item warm-up time: \(5\) seconds;
    \item measurement time: \(10\) seconds.
\end{itemize}

For close block-size comparisons, a longer controlled benchmark was used:

\begin{itemize}
    \item sample size: \(100\);
    \item warm-up time: \(10\) seconds;
    \item measurement time: \(20\) seconds.
\end{itemize}

We report the Criterion central estimate together with the measurement
interval when discussing individual optimization steps. Outliers reported
by Criterion are not silently removed from the benchmark configuration.

Because sub-millisecond timings on a laptop-class processor can move with
thermal state and background activity, differences at approximately the
one-percent level are treated cautiously and are not promoted to
algorithmic claims without controlled reruns.

\subsection{Correctness checks}

Performance measurements are accepted only after correctness tests pass.

The repository checks, among other properties:

\begin{enumerate}
    \item direct Boolean-kernel evaluation agrees with the tree evaluator;
    \item monomial-to-kernel conversion preserves the represented polynomial;
    \item parallel kernel implementations agree with the serial reference;
    \item optimized Goldilocks multiplication agrees with the reference
    multiplication;
    \item fused kernel folds agree with the unfused mathematical expression;
    \item Lagrange evaluation reconstructs the same polynomial value as the
    monomial representation on test instances.
\end{enumerate}

Optimization variants are therefore compared only after establishing that
they compute the same field value.

\subsection{Treatment of setup and precomputation}

Setup is excluded from a timed loop only when it is reusable in the
corresponding representation-level workload.

For example, converting monomial coefficients into Boolean-kernel
coefficients is not charged to the kernel single-point evaluator because
the benchmark assumes that the polynomial is already represented in that
basis. Likewise, construction of a reusable roots-of-unity domain should
not be charged repeatedly to every Lagrange query if the same domain is
reused by the prover.

Conversely, temporary work intrinsic to one evaluation belongs in the
timed path unless it is explicitly identified as reusable.

All exclusions must therefore be interpreted as representation assumptions,
not as claims that conversion or preprocessing is free.

\subsection{Lagrange comparison policy}

The repository contains both serial and parallel arbitrary-point Lagrange
evaluation over a roots-of-unity domain using batch inversion.

However, the numerical Lagrange results are not frozen at this stage of the
paper. The Boolean-kernel arithmetic backend changed substantially during
optimization, so an old Lagrange measurement collected under an earlier
field implementation is not a valid final comparison.

Before a Lagrange number is placed in the final results table, the Lagrange
implementation will be rerun using the same final field backend, machine,
compiler profile, and Criterion methodology.

This prevents the paper from combining timings collected under materially
different arithmetic implementations.

\subsection{Fairness rules}

The experimental comparison follows the following rules:

\begin{enumerate}
    \item the same field implementation is used by all bases in a final
    comparison;
    \item every method receives the same polynomial size \(N\);
    \item all timings used in a direct ratio are collected on the same
    benchmark machine;
    \item serial and parallel algorithms are labeled separately;
    \item legitimate parallelism is allowed for competing bases;
    \item representation conversion is reported separately from evaluation;
    \item reusable preprocessing is excluded only when the corresponding
    workflow can genuinely reuse it;
    \item all optimized implementations must agree with a correctness
    reference before their timings are accepted;
    \item results within the observed noise floor are not treated as
    meaningful improvements.
\end{enumerate}

These rules are designed to make the benchmark a comparison of useful
algorithms rather than of deliberately favorable implementation choices.

\section{Experimental Results}
\label{sec:results}

This section reports the final cross-basis benchmark collected in a single
run using the same machine, compiler configuration, Goldilocks backend, and
Criterion methodology for every representation.

The final comparison includes:

\begin{enumerate}
    \item the optimized Boolean-kernel evaluator;
    \item the blocked parallel monomial evaluator with block size \(4096\);
    \item serial Horner evaluation;
    \item parallel roots-of-unity Lagrange evaluation using batch inversion;
    \item serial roots-of-unity Lagrange evaluation using batch inversion.
\end{enumerate}

For every input size, all five methods represent the same deterministic
polynomial and are checked for equality at the evaluation point before
timing.

\subsection{Cross-basis scaling}

\cref{tab:final-basis-results} reports the Criterion central estimates for
all six tested sizes.

\begin{table}[ht]
\centering
\small
\begin{tabular}{rrrrrr}
\toprule
\(N\) &
Kernel &
Monomial par. &
Horner &
Lagrange par. &
Lagrange ser. \\
\midrule
\(2^{10}\) &
\(4.343\,\mu\mathrm{s}\) &
\(8.613\,\mu\mathrm{s}\) &
\(8.708\,\mu\mathrm{s}\) &
\(109.710\,\mu\mathrm{s}\) &
\(21.612\,\mu\mathrm{s}\) \\

\(2^{12}\) &
\(16.975\,\mu\mathrm{s}\) &
\(34.599\,\mu\mathrm{s}\) &
\(35.491\,\mu\mathrm{s}\) &
\(250.220\,\mu\mathrm{s}\) &
\(84.269\,\mu\mathrm{s}\) \\

\(2^{14}\) &
\(59.153\,\mu\mathrm{s}\) &
\(69.174\,\mu\mathrm{s}\) &
\(146.720\,\mu\mathrm{s}\) &
\(439.930\,\mu\mathrm{s}\) &
\(334.240\,\mu\mathrm{s}\) \\

\(2^{16}\) &
\(91.116\,\mu\mathrm{s}\) &
\(139.480\,\mu\mathrm{s}\) &
\(558.910\,\mu\mathrm{s}\) &
\(1.2503\,\mathrm{ms}\) &
\(1.3391\,\mathrm{ms}\) \\

\(2^{18}\) &
\(263.230\,\mu\mathrm{s}\) &
\(380.910\,\mu\mathrm{s}\) &
\(2.2440\,\mathrm{ms}\) &
\(4.5612\,\mathrm{ms}\) &
\(5.3679\,\mathrm{ms}\) \\

\(2^{20}\) &
\(1.0560\,\mathrm{ms}\) &
\(1.4116\,\mathrm{ms}\) &
\(9.0402\,\mathrm{ms}\) &
\(17.1990\,\mathrm{ms}\) &
\(21.3890\,\mathrm{ms}\) \\
\bottomrule
\end{tabular}
\caption{Final same-run arbitrary-point evaluation benchmark over the
Goldilocks field. Lower is better.}
\label{tab:final-basis-results}
\end{table}

The Boolean-kernel evaluator is the fastest method at every tested size in
this final sweep.

\subsection{Headline result at \(N=2^{20}\)}

For
\[
N=2^{20}=1{,}048{,}576,
\]
the optimized Boolean-kernel evaluator has Criterion interval
\[
[1.0009,\ 1.1132]\ \mathrm{ms}
\]
with central estimate
\[
\boxed{1.0560\ \mathrm{ms}}.
\]

The corresponding throughput is approximately
\[
\boxed{9.93\times10^8}
\]
input coefficients per second.

The parallel monomial baseline has central estimate
\[
1.4116\ \mathrm{ms}.
\]
Thus, in the final same-run comparison,
\[
\frac{1.4116}{1.0560}
=
\boxed{1.337\times}.
\]

Equivalently, the Boolean-kernel evaluator uses approximately
\[
1-\frac{1.0560}{1.4116}
\approx
\boxed{25.2\%}
\]
less wall-clock time than the blocked parallel monomial evaluator.

Against serial Horner evaluation,
\[
\frac{9.0402}{1.0560}
=
\boxed{8.56\times}.
\]

\subsection{Comparison with the Lagrange representation}

At
\[
N=2^{20},
\]
parallel arbitrary-point evaluation from the roots-of-unity Lagrange
representation requires
\[
17.199\ \mathrm{ms},
\]
while the serial version requires
\[
21.389\ \mathrm{ms}.
\]

Relative to the optimized Boolean-kernel evaluator, these correspond to
\[
\frac{17.199}{1.0560}
=
\boxed{16.29\times}
\]
and
\[
\frac{21.389}{1.0560}
=
\boxed{20.25\times},
\]
respectively.

This result concerns \emph{arbitrary-point evaluation from the stored
representation}.  It should not be interpreted as evidence that
roots-of-unity evaluation form is generally inferior in a proof system.
Evaluation form is highly advantageous for other workloads, including
pointwise operations and FFT/NTT-based encoding.  The comparison instead
shows that an arbitrary off-domain evaluation has a very different cost
profile in the three representations considered here.

\subsection{Scaling relative to the kernel evaluator}

The ratio between each competitor and the Boolean-kernel evaluator changes
with \(N\).  For the fair parallel monomial baseline, the measured ratios
are approximately

\[
\begin{array}{c|cccccc}
N &
2^{10} &
2^{12} &
2^{14} &
2^{16} &
2^{18} &
2^{20}
\\
\hline
T_{\mathrm{mono,par}}/T_{\mathrm{kernel}}
&
1.98 &
2.04 &
1.17 &
1.53 &
1.45 &
1.34
\end{array}
\]

The non-monotonic behavior reflects the interaction between subtree
parallelism, Rayon scheduling, cache behavior, and the fixed block size.
It is therefore more appropriate to report the measured curve than to
infer a constant hardware-independent speedup factor.

For serial Horner evaluation, the ratio grows much more strongly with
problem size:
\[
2.00,\ 2.09,\ 2.48,\ 6.13,\ 8.52,\ 8.56.
\]
This is consistent with the fact that the kernel tree exposes substantial
intra-polynomial parallelism while Horner evaluation retains a long
sequential dependency chain.

\subsection{Best observed optimization run versus final comparison}

During optimization of the Boolean-kernel evaluator alone, the best
observed \(N=2^{20}\) run was
\[
\boxed{892.57\ \mu\mathrm{s}},
\]
with interval
\[
[884.25,\ 902.65]\ \mu\mathrm{s}.
\]

The final multi-method sweep measured the same optimized kernel evaluator at
\[
1.0560\ \mathrm{ms}.
\]

We deliberately use the latter number for the cross-basis headline ratio,
because it was collected in the same benchmark campaign as all competing
representations.  The lower \(892.57\,\mu\mathrm{s}\) measurement is retained
as an optimization result and evidence of run-to-run variability, but it is
not mixed with competitor timings from a different run to construct a
speedup claim.

This distinction prevents accidental cherry-picking.

\subsection{Optimization ablation}

The optimization campaign reduced the best observed subtree-parallel runtime
approximately as follows:
\[
1.2530
\rightarrow
1.2007
\rightarrow
1.1619
\rightarrow
1.0179
\rightarrow
0.93533
\rightarrow
0.89257
\quad\mathrm{ms}.
\]

The final value is approximately
\[
28.8\%
\]
below the earlier \(1.2530\)-ms subtree baseline.

The major experimentally validated improvements were:

\begin{enumerate}
    \item specialized Goldilocks multiplication reduction;
    \item fused field-level kernel fold;
    \item elimination of bounds checks in the hot level loop;
    \item direct first-level reduction into half-size scratch;
    \item worker-local scratch reuse.
\end{enumerate}

The isolated microbenchmarks showed approximately:

\begin{center}
\begin{tabular}{lrrr}
\toprule
Optimization &
Before &
After &
Local improvement \\
\midrule
Goldilocks multiplication &
\(2.2372\) ns &
\(2.0559\) ns &
\(8.1\%\) \\

Fused kernel fold &
\(4.0591\) ns &
\(3.3585\) ns &
\(17.3\%\) \\

Unchecked level loop &
\(6.9671\,\mu\mathrm{s}\) &
\(6.3575\,\mu\mathrm{s}\) &
\(8.75\%\) \\

Direct first level &
\(6.9031\,\mu\mathrm{s}\) &
\(6.1571\,\mu\mathrm{s}\) &
\(10.8\%\) \\
\bottomrule
\end{tabular}
\end{center}

These percentages describe the affected primitive or stage and should not be
added together to predict whole-evaluator speedup.

\subsection{Negative results}

Two-level and three-level tree fusion were both rejected.  In particular,
the two-level experiment gave
\[
9.7019\ \mu\mathrm{s}
\]
for the ordinary two-level path and
\[
9.8528\ \mu\mathrm{s}
\]
for the fused path.

Thus the version with fewer intermediate memory accesses was approximately
\(1.6\%\) slower.

This result illustrates an important implementation constraint: reducing
memory traffic does not necessarily improve scalar field code when doing so
increases dependency depth and register pressure.

\subsection{Interpretation}

The final results establish a narrow but concrete claim:

\begin{quote}
For arbitrary single-point evaluation of a degree-\(<N\) univariate
polynomial already stored in the tested representation, the optimized
Boolean-kernel evaluator is faster than the monomial and roots-of-unity
Lagrange implementations in our Rust benchmark over Goldilocks on the
tested Apple M1 machine.
\end{quote}

At \(N=2^{20}\), the conservative same-run comparison gives a
\[
1.337\times
\]
advantage over the parallel monomial baseline and a
\[
25.2\%
\]
reduction in wall-clock evaluation time.

The experiment does not yet establish an end-to-end prover speedup.  That
requires measuring basis conversion, encoding, commitment, FRI-related
operations, hashing, and all other protocol costs in an integrated system.

\section{Relevance to Proof Systems and Limitations}
\label{sec:proof-systems}

The experiments in this paper isolate arbitrary single-point evaluation of a
univariate polynomial already stored in a chosen representation.  This
section explains where such a primitive appears in existing proof-system
architectures, which parts of the optimization can transfer directly, and
which conclusions would require end-to-end integration experiments.

\subsection{Where arbitrary-point evaluations appear}

Transparent proof systems based on Reed--Solomon proximity testing commonly
move between polynomial coefficients and evaluations on structured domains.
They also evaluate committed polynomials at verifier-derived challenge
points that need not lie on the original commitment domain.

RISC Zero provides a concrete public example of this pattern.  Its STARK
documentation describes conversion of trace columns to polynomial
coefficients with an inverse NTT, Reed--Solomon evaluation on an expanded
domain, DEEP-ALI checks at an out-of-domain point \(z\), and a subsequent
FRI low-degree argument \cite{risc0starkbyhand,risc0sequence}.

In the DEEP-ALI stage, the prover supplies evaluations of trace polynomials
at points of the form
\[
P_i(\omega^n z),
\]
where \(z\) is a verifier challenge and the shifts depend on the constraint
system \cite{risc0sequence}.  This is the class of operation for which an
efficient arbitrary-point evaluator can potentially matter.

The relevance is nevertheless conditional on representation.  Existing
STARK implementations are typically organized around evaluation-domain
vectors because FFT/NTT-based encoding, pointwise constraint operations,
Merkle commitments, and FRI naturally consume such data.  Replacing only
one evaluation routine without considering the representation of the
surrounding pipeline may therefore introduce conversion costs that dominate
the local gain.

\subsection{Relation to FRI}

FRI proves that a committed vector is close to evaluations of a low-degree
polynomial.  Its commit phase recursively folds evaluations into smaller
commitments, while the query phase opens a small number of positions from
the committed layers \cite{risc0fri}.

The Boolean-kernel fold studied in this paper,
\[
F_t(a,b)=t(a+b)+b,
\]
should not be identified with a standard FRI fold merely because both are
binary reductions.  They act on different representations and satisfy
different protocol roles.

The present result instead motivates the following systems question:

\begin{quote}
Can a prover retain a Boolean-kernel representation across enough of the
polynomial pipeline that the fast arbitrary-point evaluator is useful
without paying an offsetting conversion cost before FRI or commitment?
\end{quote}

Answering this requires protocol-level design and benchmarking.  The current
paper establishes the performance of the local representation primitive;
it does not establish a new FRI protocol.

\subsection{Transferable low-level optimizations}

Some optimizations presented in \cref{sec:implementation} are not specific
to the Boolean-kernel basis.

The specialized Goldilocks reduction exploits only
\[
p=2^{64}-2^{32}+1.
\]
It can therefore benefit other Goldilocks workloads when the same arithmetic
pattern occurs.

This is relevant because Goldilocks has already been used specifically to
obtain efficient CPU arithmetic in FRI-based proof systems.  Plonky2, for
example, was designed around FRI and the Goldilocks field precisely because
the modulus permits efficient \(64\)-bit arithmetic on modern processors
\cite{plonky2intro}.

Likewise, the experiments on allocation reuse, worker-local scratch space,
bounds-check elimination, and task granularity are general implementation
lessons for prover kernels.  Their exact gains are workload-dependent, but
they do not rely on a change to the cryptographic protocol.

By contrast, the speedup arising from the Boolean-kernel representation
itself transfers only if the relevant polynomial is actually available in
that representation.

\subsection{Effect on an existing prover}

Suppose an existing prover currently represents a polynomial by monomial
coefficients or by evaluations on a roots-of-unity domain.

There are three distinct integration cases.

\paragraph{Case 1: local arithmetic replacement only.}
If the prover retains its existing representation and proof format but
adopts only the Goldilocks arithmetic optimizations, then the cryptographic
protocol is unchanged.  Any speedup is limited to the affected arithmetic
kernels.

\paragraph{Case 2: internal representation change with identical proof
interface.}
If the prover internally stores some polynomials in Boolean-kernel form but
converts back before commitment or proof serialization, then an external
verifier need not change.  The benefit is positive only if
\[
T_{\mathrm{saved\ evaluations}}
>
T_{\mathrm{basis\ conversions}}
+
T_{\mathrm{integration\ overhead}}.
\]

\paragraph{Case 3: protocol-visible representation change.}
If commitments, openings, FRI layers, transcript messages, or verifier
equations themselves change, then the proof format changes.  Existing
verifiers cannot automatically accept the new proof format.

This distinction is important when discussing systems that already have
deployed verification infrastructure.

\subsection{On-chain verification}

A prover optimization does not by itself reduce on-chain verification gas.

For example, RISC Zero publishes Solidity verifier interfaces for receipts
and Ethereum integration contracts \cite{risc0ethverifier}.  These contracts
consume a particular proof or receipt format.

If a faster prover produces exactly the same proof bytes and verification
relation, then:
\[
\boxed{
\text{prover time may decrease while on-chain verification remains unchanged}.
}
\]

In this case no smart-contract change is required, but gas usage is not
expected to improve merely because the prover computed the proof faster.

If instead the Boolean-kernel representation changes the proof object or
verification equations, then deployed contracts would need a compatible new
verifier implementation, a routing/versioning mechanism, or a migration
path.  Existing RISC Zero Ethereum infrastructure, for example, uses
verifier selectors and routing between verifier implementations, illustrating
that verifier parameters and proof formats are versioned objects
\cite{risc0router}.

Accordingly, the present benchmark should be viewed primarily as a
\emph{prover-side} result.

\subsection{When the benchmark is likely to matter}

The representation is most promising in workloads satisfying several of
the following conditions:

\begin{enumerate}
    \item large polynomials are evaluated repeatedly at verifier-derived
    off-domain points;
    \item the polynomial can remain in Boolean-kernel representation for
    more than one operation;
    \item basis conversion can be amortized or eliminated;
    \item the prover is CPU-bound in field arithmetic or memory movement
    rather than dominated by hashing or external I/O;
    \item the target field admits efficient realization of the fold
    arithmetic.
\end{enumerate}

Conversely, the result is less likely to affect workloads in which the
polynomial is used almost exclusively through evaluations on one fixed FFT
domain or where a single arbitrary-point evaluation is surrounded by much
larger hashing, commitment, or encoding costs.

\subsection{Representation conversion}

The monomial-to-kernel conversion described in
\cref{sec:representations} is a Boolean Möbius transform and costs
\[
O(N\log N)
\]
field operations.

This cost is not included in the single-point evaluation timing because
the benchmark assumes the input is already represented in the basis being
measured.

For a one-off query, conversion can easily outweigh the evaluation gain.
The basis is therefore most interesting when either:

\begin{enumerate}
    \item the witness or polynomial is generated directly in kernel
    coefficients;
    \item several subsequent operations consume the kernel representation;
    \item conversion is fused with another prover stage;
    \item or the representation replaces another transform already present
    in the prover.
\end{enumerate}

An end-to-end evaluation must account for the full representation lifetime,
not only the isolated query.

\subsection{Limitations of the present experiments}

The current study has several explicit limitations.

\paragraph{Single processor family.}
The headline measurements are from an Apple M1 system.  Cache sizes,
instruction scheduling, multiplication latency, memory bandwidth, and the
cost of Rayon scheduling differ on x86-64 servers and other ARM systems.

\paragraph{Single base field.}
The final measurements use Goldilocks.  The Boolean-kernel algebra is
field-independent, but the constant factors of the implementation are not.
A field such as Mersenne-31 may produce a different balance between
arithmetic and memory costs.

\paragraph{CPU implementation.}
No GPU, FPGA, or wide-SIMD implementation is evaluated in this paper.

\paragraph{Single-point evaluation focus.}
The paper does not yet benchmark all operations needed by a proof system.
In particular, full-domain encoding, low-degree extension, commitment,
basis conversion, and FRI integration are separate workloads.

\paragraph{No end-to-end prover measurement.}
The measured
\[
1.337\times
\]
same-run speed advantage over the parallel monomial baseline applies to the
single-point evaluator only.  It is not an end-to-end STARK or PCS speedup.

\paragraph{Run-to-run variability.}
The best kernel-only optimization run reached
\[
892.57\,\mu\mathrm{s},
\]
while the final same-run cross-basis sweep measured
\[
1.0560\,\mathrm{ms}.
\]
The paper therefore uses same-run measurements when constructing
cross-basis ratios and reports the lower number only as a best observed
implementation result.

\paragraph{Lagrange implementation scope.}
The roots-of-unity Lagrange benchmark uses the direct barycentric formula
with batch inversion for an arbitrary point.  It is not intended to replace
FFT/NTT evaluation when the requested points form a structured domain.

\subsection{Next system-level experiments}

The benchmark naturally suggests the following next experiments:

\begin{enumerate}
    \item measure monomial-to-kernel and kernel-to-monomial conversion
    alongside single-point evaluation;
    \item measure repeated queries to determine the query count at which
    conversion amortizes;
    \item integrate the evaluator into a DEEP-style off-domain evaluation
    stage of an existing STARK implementation;
    \item benchmark a kernel-native route to FRI commitment and folding;
    \item repeat the study on x86-64, Mersenne-31, and additional processor
    architectures;
    \item measure the fraction of total prover time attributable to the
    operation before making an end-to-end performance claim.
\end{enumerate}

The principal open engineering question is therefore no longer whether the
Boolean-kernel representation supports a fast arbitrary-point evaluation.
The experiments establish that it does.  The remaining question is whether
a complete proof-system pipeline can preserve the representation long enough
for this local advantage to survive at system scale.

\section{Conclusion}
\label{sec:conclusion}

We studied arbitrary single-point evaluation of degree-\(<N\) univariate
polynomials represented in the Boolean-kernel basis
\[
K_y(X)
=
\prod_{i=0}^{m-1}
\left(X^{2^i}+y_i\right),
\qquad
N=2^m,
\qquad
y\in\{0,1\}^m.
\]

The representation admits the local fold
\[
F_t(a,b)=t(a+b)+b,
\qquad
t=x^{2^i},
\]
which evaluates the polynomial through a complete binary reduction tree.
The tree contains exactly
\[
N-1
\]
folds and therefore requires exactly
\[
N-1
\]
field multiplications and
\[
2(N-1)
\]
field additions, together with
\[
m-1
\]
squarings to generate the level parameters.

On the implementation side, the final Rust evaluator combines
coarse-grained subtree parallelism, specialized Goldilocks reduction,
a fused kernel fold, bounds-check elimination in the hot loop, direct
first-level reduction, and worker-local scratch reuse.

For
\[
N=2^{20},
\]
the final same-run benchmark on the tested Apple M1 system measured
\[
1.0560\ \mathrm{ms}
\]
for the optimized Boolean-kernel evaluator and
\[
1.4116\ \mathrm{ms}
\]
for the blocked parallel monomial baseline.  This corresponds to a
\[
\boxed{1.337\times}
\]
evaluation-speed advantage, or approximately
\[
\boxed{25.2\%}
\]
lower wall-clock time for the isolated arbitrary-point evaluation
primitive.

A separate kernel-only optimization run reached
\[
892.57\,\mu\mathrm{s},
\]
but we do not combine that number with competitor timings collected in a
different run when computing cross-basis speedups.

The comparison with the roots-of-unity Lagrange representation shows an
even larger difference for arbitrary off-domain evaluation.  At
\(N=2^{20}\), the measured parallel Lagrange evaluator required
\(17.199\) ms.  This does not imply that evaluation form is generally
inferior in proof systems; rather, it emphasizes that representation choice
is workload-dependent.

The central systems question is therefore not whether the Boolean-kernel
basis supports fast single-point evaluation---the experiments establish that
it does---but whether a prover can retain this representation across enough
of its pipeline that conversion, encoding, commitment, and FRI-related
costs do not erase the local advantage.

Future work should measure that representation lifetime directly inside a
complete transparent proof system, including basis conversion, repeated
off-domain queries, commitment, and FRI integration.  Only such an
end-to-end study can determine how much of the isolated evaluator speedup
survives at prover scale.


\bibliographystyle{alpha}
\bibliography{references}

\end{document}
